III.34
pratibha va sarvam

III.35
hrdaye citta-samvit

III.36
sattva-purusha atyanta-asankirna pratyaya-avisesa bhoga para-arthat svartha-samyamat purusha-jnana

III.37
tata pratibha-sravana-vedana-adarsa-asvada-vartta jayante

III.38
te samadhav upasarga vyutthane siddhaya


III.34
pratibha va sarvam

A cognition is either a concept or an intuition.
The above sutra marks a very important division in Jnana-pada.
I divide Yoga into prajna-kriya-jnana-dharma padas.
I also divide this into prajna kriya-jnana dharma in that
kriya-jnana is six-fold sankhya yoga and needs to be considered
"as one" traya ekatra samyama. three as one. actually this is 
inner 3-as-1. Six-fold includes outer 3-as-1...pratyahara
This is important because only inner 3-as-1 is
raised to pure reason, not outer 3-as-1.

Sorry to go off but I need to show a connection 
between the VA in prajna-pada and the VA in jnana-pada
samyama is jnana from prajna. I need to ask some assistance.
I am trying to present a pretty radical criticism of Vedanta-Yoga.

Here's the point. A cognition is either a concept or an intuition.
This is central in that the VA in jnana-pada really is top-level
division and in Logic a top-level division, a Root Divide.
....is key.....and I am saying Vedantins are missing this key.

The key is that Vedanta is practicing an Intuitive Yoga.
Its not hard to see a connection between those five sutras.
the "tata" shows that pratibha is the fruit of purusha-jnana.
So purusha-jnana is what Vedanta Yoga is all about.
Now Patanjali says Intuitive Purusha Jnana Yoga is 
an obstacle to samadhi. That is key.
The five samyama bhumis starting here are purusha bhumi.
The previous five samyasma bhumis are prakriti bhumis.

So we need to divide the 11 samyama bhumis into a triad
because a pentad indicates a constructed unity as a monad
This division is Prakriti-Purusha-Buddhi
This is important to the concept of samyoga.
You have to have mastered Samkhya to get this.
Vedanta and Samkhya are both prerequisites of Yoga.

Vedanta is not really practicing conceptual Yoga.
And by concept I mean ordinary science.
The fact that they call Jnanapada Vibhutipada definitely
should tell you something. They are advertising that they
are all about Spiritual Powers. The Path of Seeing is Purusha-jnana
and that is eyes-closed videha yoga .
This whole path and this is just the start.
The heart of this path is Mahavideha Samyama.

This is still Sabija Yoga because it is based on
Skill in Nirvicara. 

Knowledge born from discrimination is what 
fills the gap between Knowledge born from Nature
and knowledge born from Spirit.
[sarva visaya vs sarvata visaya]
But knowledge born from discrimination is pure reason.
Buddhi in itself. It is the first objectification
of Absolute Reason. Buddhi is a tattva and is the
first objectification of Reason. It cant be
known from ordinary samadhi since it is the specification of
ordinary samadhi, the Idea of Samadhi. kshana.tat.krama.
and that is the absolute idea. the middle way between 
kshana and krama